\section{Scope and Evidence Method}
\label{sec:scope}
\paragraph{Unit of analysis.}
We consider research on compositional layouts, design-specific visual text, vector or layered artifact construction, document editing and reconstruction, and design-specific evaluation.
The intended domain is defined by the communication and composition task, rather than the appearance of a design-themed image.
We retain generic image-generation work only when an explicit graphic-design constraint, component structure, or evaluation is central.
Table~\ref{tab:task-grounding} links representative application families to original studies and their actual input/output contracts.

% Independent source table.
\begin{table*}[t]
\centering
\caption{Concrete graphic-design research tasks and representative primary studies.
The input/output description is a task-level summary; it should not be interpreted as a guarantee of native artifact editability across all cited papers.}
\label{tab:task-grounding}
\footnotesize
\begin{tabularx}{\textwidth}{@{}p{0.21\textwidth}p{0.34\textwidth}X@{}}
\toprule
\textbf{Application / research task} & \textbf{Representative primary work} & \textbf{Input, output and evaluation boundary}\\
\midrule
Poster layout and creation & \ac{posterlayout} \citep{posterlayout2023}; \ac{vinci} \citep{vinci2021}; \ac{posteragent} \citep{posteragent2026} & Product/image canvas or design brief $\rightarrow$ positioned visual/text elements or complete poster; separate layout from final file structure.\\
Advertising and marketing design & \ac{cglgan} \citep{cglgan2022}; \ac{layoutdetr} \citep{layoutdetr2024}; \ac{autoposter} \citep{autoposter2023} & Product/background assets $\rightarrow$ advertising composition; content-aware placement and full creative rendering are distinct tasks.\\
Presentation slide generation & \ac{autopresent} \citep{autopresent2025}; \ac{pptagent} \citep{pptagent2025} & Natural-language request or source document $\rightarrow$ slide code/actions or presentation; distinguish single slides from slide-deck coherence.\\
Infographics and information graphics & \ac{bizgen} \citep{bizgen2025}; \ac{infoedit} \citep{yang2026infoedit} & Article $\rightarrow$ visually dense infographic, or existing infographic + instruction $\rightarrow$ edited result requiring reflow.\\
Scientific illustration / academic posters & \ac{autofigure} \citep{autofigure2026}; \ac{postergen} \citep{postergen2026} & Scientific text or paper $\rightarrow$ illustration or conference poster; factual content, spatial hierarchy and output format all matter.\\
Typography and vector icons & \ac{textdiffuser} \citep{textdiffuser2023}; \ac{wordasimage} \citep{wordasimage2023}; \ac{iconshop} \citep{iconshop2023} & Glyph/text prompt $\rightarrow$ rendered letterforms or vector paths; modifying glyph contours is not the same as editing semantic text.\\
\bottomrule
\end{tabularx}
\end{table*}


\paragraph{Discovery and selection.}
The working evidence register is seeded from a public research catalogue, prior surveys, original research proceedings, author project pages, and references within candidate papers.
A publication and its preprint versions should be deduplicated as one research contribution after verifying their relationship.
The catalogue is a living research index, not a final screened sample.
We code multi-task papers on multiple analytical axes, even if they have only one primary catalogue category.
Benchmark selection includes \ac{rico}, \ac{publaynet}, and \ac{doclaynet} as common layout and structured-document contexts, while distinguishing their domain and label differences rather than pooling results.

\paragraph{P1 original-source citation expansion.}
Ten original-source-identified P1 studies expand the bibliography to 196 distinct works, linking 167 of the 225 working register IDs and leaving 58 pending identity or scope review.
A subsequent single-study original-methods/results audit brings the bibliography to 197 works and links 168 working-study IDs, with 57 records still awaiting identity or scope review.
This audit separately examines layer-reconstruction metrics and qualitative editing examples instead of treating them as the same evaluation.
The newly linked references remain abstract/introduction-level citation audits, not independently verified experiments.

\paragraph{Verification levels.}
We distinguish \emph{catalogue-only} entries, \emph{author/publisher abstract or indexed} entries, and works for which at least selected \emph{original methods or results} have been inspected.
The last group is not necessarily a full-text audit of all appendices or an independent replication.
The manuscript also maintains a citation-level audit register for all bibliography entries, but the new registry starts with unverified source-depth fields rather than transferring aggregate verification counts to unlinked citation keys.
For each cited work, source-level verification must be linked to a primary-paper section or figure and reconciled against the broader study register before an evidence tier is promoted.
Empirical statements require identifiable evidence such as a method section, table, figure, or publisher record.
Missing evaluation details in an uninspected paper are not treated as evidence that no evaluation exists.

At the November 3, 2026 audit checkpoint, the broader provisional register held 225 entries: 88 with selected original method/result sections checked, 28 supported through abstract or publication-index evidence, and 109 catalogue-only.
These figures describe \emph{verification depth}, not a \ac{prisma} screening outcome, a claim that all 225 studies were comprehensively read, or the size of the stable public catalogue.
At this checkpoint, the bibliography contains 180 cited references, including 36 additions whose original abstract and bibliographic identities were checked, but whose experiment details have not all been inspected.
A conservative identifier/title crosswalk links 138 of the 207 public catalogue papers to bibliography entries and flags 69 unresolved cases for further adjudication.
The initial frozen spreadsheet crosswalk linked 148 working-study IDs to 180 cited bibliography entries and left 77 records for identity review.
After confirming three publication/preprint version identities and introducing six original-source-checked studies, the 186-entry bibliography now links 157 of the 225 working-study IDs and leaves 68 pending identity or coverage review; these are not exclusions.

\paragraph{Coding protocol.}
For each work, we record input and task, output representation, control mechanism, editing operation, measurement target, metric/protocol, and source depth.
Boundaries matter: \ac{svg} paths are not necessarily editable semantic text; raster inpainting is a real image-editing capability without implying Photoshop-style object layers; and a graphical user interface can be demonstrated without a controlled user study.

\paragraph{Limitations of the review process.}
This is an initial narrative survey with structured coding, \emph{not yet} a reproducible systematic review.
Search strings, database query dates, deduplication logs, exclusion rules, independent screening decisions, and reviewer agreement must be finalized before presenting complete coverage.
The bibliography is deliberately representative rather than a claim to cite every study in the current registry.
In the remainder, we avoid pooling incomparable metrics into a single leaderboard.
